% Volume~II, abstracts (writer "results"): the thesis's own language first, then the other,
% as in Volume~I. The master file calls it after the title page (% Volume~II, abstracts (writer "results"): the thesis's own language first, then the other,
% as in Volume~I. The master file calls it after the title page (% Volume~II, abstracts (writer "results"): the thesis's own language first, then the other,
% as in Volume~I. The master file calls it after the title page (% Volume~II, abstracts (writer "results"): the thesis's own language first, then the other,
% as in Volume~I. The master file calls it after the title page (\input{resume}).

\chapter*{Abstract}
\addcontentsline{toc}{chapter}{Abstract}
\markboth{Abstract}{Abstract}
\begin{otherlanguage}{english}
Volume~I of \emph{troncons} measured an ordinary laptop (16\,GB of RAM, an 8\,GB RTX\,4060, an encrypted NVMe disk), proved the roof of the problem and decided, by calculation, to build an exact architecture first: serve through the graphics card a model much larger than all the memory of the machine, its experts compressed losslessly on disk and decoded at the moment a computation needs them, within a memory window that does not depend on the size of the model. This second volume is the volume of the system: the program, built in Rust with CUDA kernels by a team of five lots with frozen interfaces, and its first milestone, judged on gpt-oss-120b (63.4\,GB, four times the RAM, eight times the card's memory) against twelve criteria declared before the first line of code.

The volume gives the architecture and structure of the program; the \texttt{.trc} format to the bit, one chunk per expert and streams stored by refinement level; the byte path, from the disk to the card's memory, through unbuffered reads and pinned memory; the kernels, among them the rANS decoder with per-lane tables, with the proof of their exactness and determinism; a double-precision reference and the theory of the judge, founded on the Kullback--Leibler divergence; the expert cache, the scheduler and the engine, with their lemmas.

The results fit in four numbers, each recomputed by a second path from the committed passes. \emph{Memory}: a 384\,MB peak for the 63.4\,GB model, the same to within 3.7\,MB as for a model five times smaller, where llama.cpp holds 10.9\,GB (11.4\,GB in its fastest configuration); memory no longer follows the size of the model. \emph{Exactness}: logits identical to the bit between the coded and the raw file over 5\,120 steps, from one build and one order of computation to another. \emph{Speed}: 1.965 tokens per second with the integrity guard, 2.22 times the best configuration of llama.cpp, and a prefill seven times faster; two earlier ratios, $\times 3.85$ then $\times 2.79$, are struck out because the baseline was badly tuned (sweeping it raised it from 0.57 to 0.88 token/s). \emph{Correctness}: closer to the exact mathematics of the model than llama.cpp (KL 0.0031 against 0.0607 on the 120b), and, with an f32 KV cache, equal to them up to the format: the whole gap came from the f16 rounding of the cache.

A judge who wrote neither the specification nor the code redid the numbers: his decoder, written from the specification alone, reproduces the chunks to the bit; his own double-precision forward pass confirms the reference up to the rounding of the format. He found two serious defects (a pass that hangs without a message, an orphan child process that blocks its driver) and five medium ones, among them a corruption of raw bytes invisible down to the text. The round of corrections corrected or measured them in one morning, each with a control that fails on the old binary; a production guard now turns any corruption of a chunk read into a named fatal error, for 1.6\,\% of speed at equal disk. Ten criteria out of twelve hold; the other two await a decision of the arbiter.

What remains is the research contribution, the base + complement split of milestone J2; an unencrypted disk, which will tell what BitLocker costs on a real model; Stone's biomimetic study, the weight genome; and the gate that will give this engine back to relais. The thesis exists in two complete versions, English and French; every figure is regenerated from the committed measurement data.
\end{otherlanguage}

\chapter*{Résumé}
\addcontentsline{toc}{chapter}{Résumé}
\markboth{Résumé}{Résumé}
\begin{otherlanguage}{french}
Le Tome~I de \emph{troncons} avait mesuré un portable ordinaire (16\,Go de RAM, une RTX\,4060 de 8\,Go, un disque NVMe chiffré), démontré le toit du problème et décidé, par le calcul, de construire d'abord une architecture exacte : servir par la carte graphique un modèle bien plus gros que toute la mémoire de la machine, ses experts compressés sans perte sur le disque et décodés à l'instant où un calcul en a besoin, dans une fenêtre de mémoire qui ne dépend pas de la taille du modèle. Ce second tome est celui du système : le programme, construit en Rust avec des noyaux CUDA par une équipe de cinq lots aux interfaces gelées, et son premier jalon, jugé sur gpt-oss-120b (63,4\,Go, quatre fois la RAM, huit fois la mémoire de la carte) contre douze critères déclarés avant la première ligne de code.

Le tome donne l'architecture et la structure du programme ; le format \texttt{.trc} au bit près, un tronçon par expert et des flux rangés par niveau de raffinement ; le chemin des octets, du disque à la mémoire de la carte, par la lecture sans tampon et la mémoire épinglée ; les noyaux, dont le décodeur rANS à table par voie, avec la preuve de leur exactitude et de leur déterminisme ; une référence en double précision et la théorie du juge, fondée sur la divergence de Kullback et Leibler ; le cache d'experts, l'ordonnanceur et le moteur, avec leurs lemmes.

Les résultats tiennent en quatre nombres, chacun recalculé par un second chemin depuis les passes commitées. \emph{Mémoire} : 384\,Mo de pointe pour le modèle de 63,4\,Go, la même à 3,7\,Mo près que pour un modèle cinq fois plus petit, quand llama.cpp en tient 10,9\,Go (11,4\,Go dans sa configuration la plus rapide) ; la mémoire ne suit plus la taille du modèle. \emph{Exactitude} : des logits identiques au bit entre le fichier codé et le fichier brut sur 5\,120 pas, d'une construction et d'un ordre de calcul à l'autre. \emph{Vitesse} : 1,965 jeton par seconde avec la garde d'intégrité, 2,22 fois la meilleure configuration de llama.cpp, et un préremplissage sept fois plus rapide ; deux rapports antérieurs, $\times 3{,}85$ puis $\times 2{,}79$, sont barrés parce que la base était mal réglée (la balayer l'a portée de 0,57 à 0,88 jeton/s). \emph{Justesse} : plus près des mathématiques exactes du modèle que llama.cpp (KL 0,0031 contre 0,0607 sur le 120b), et, avec un cache KV en f32, égal à elles au format près : tout l'écart venait de l'arrondi f16 du cache.

Un juge qui n'avait écrit ni la spécification ni le code a refait les chiffres : son décodeur, écrit depuis la seule spécification, redonne les tronçons au bit près ; sa propre passe avant en double précision confirme la référence à l'arrondi du format près. Il a trouvé deux défauts graves (une passe qui se fige sans message, un processus enfant orphelin qui bloque son pilote) et cinq moyens, dont une corruption d'octets bruts invisible jusque dans le texte. La ronde de corrections les a corrigés ou mesurés en une matinée, chacun avec un témoin qui tombe sur l'ancien binaire ; une garde de production rend désormais toute corruption d'un tronçon lu en erreur fatale nommée, pour 1,6\,\% de vitesse à disque égal. Dix critères tiennent sur douze ; les deux autres attendent une décision de l'arbitre.

Restent la contribution de recherche, la découpe base + complément du jalon J2 ; un disque non chiffré, qui dira ce que BitLocker coûte sur un vrai modèle ; l'étude biomimétique de Stone, le génome des poids ; et la porte qui rendra ce moteur à relais. Le mémoire existe en deux versions complètes, anglaise et française ; chaque figure se régénère depuis les données de mesure commitées.
\end{otherlanguage}
).

\chapter*{Abstract}
\addcontentsline{toc}{chapter}{Abstract}
\markboth{Abstract}{Abstract}
\begin{otherlanguage}{english}
Volume~I of \emph{troncons} measured an ordinary laptop (16\,GB of RAM, an 8\,GB RTX\,4060, an encrypted NVMe disk), proved the roof of the problem and decided, by calculation, to build an exact architecture first: serve through the graphics card a model much larger than all the memory of the machine, its experts compressed losslessly on disk and decoded at the moment a computation needs them, within a memory window that does not depend on the size of the model. This second volume is the volume of the system: the program, built in Rust with CUDA kernels by a team of five lots with frozen interfaces, and its first milestone, judged on gpt-oss-120b (63.4\,GB, four times the RAM, eight times the card's memory) against twelve criteria declared before the first line of code.

The volume gives the architecture and structure of the program; the \texttt{.trc} format to the bit, one chunk per expert and streams stored by refinement level; the byte path, from the disk to the card's memory, through unbuffered reads and pinned memory; the kernels, among them the rANS decoder with per-lane tables, with the proof of their exactness and determinism; a double-precision reference and the theory of the judge, founded on the Kullback--Leibler divergence; the expert cache, the scheduler and the engine, with their lemmas.

The results fit in four numbers, each recomputed by a second path from the committed passes. \emph{Memory}: a 384\,MB peak for the 63.4\,GB model, the same to within 3.7\,MB as for a model five times smaller, where llama.cpp holds 10.9\,GB (11.4\,GB in its fastest configuration); memory no longer follows the size of the model. \emph{Exactness}: logits identical to the bit between the coded and the raw file over 5\,120 steps, from one build and one order of computation to another. \emph{Speed}: 1.965 tokens per second with the integrity guard, 2.22 times the best configuration of llama.cpp, and a prefill seven times faster; two earlier ratios, $\times 3.85$ then $\times 2.79$, are struck out because the baseline was badly tuned (sweeping it raised it from 0.57 to 0.88 token/s). \emph{Correctness}: closer to the exact mathematics of the model than llama.cpp (KL 0.0031 against 0.0607 on the 120b), and, with an f32 KV cache, equal to them up to the format: the whole gap came from the f16 rounding of the cache.

A judge who wrote neither the specification nor the code redid the numbers: his decoder, written from the specification alone, reproduces the chunks to the bit; his own double-precision forward pass confirms the reference up to the rounding of the format. He found two serious defects (a pass that hangs without a message, an orphan child process that blocks its driver) and five medium ones, among them a corruption of raw bytes invisible down to the text. The round of corrections corrected or measured them in one morning, each with a control that fails on the old binary; a production guard now turns any corruption of a chunk read into a named fatal error, for 1.6\,\% of speed at equal disk. Ten criteria out of twelve hold; the other two await a decision of the arbiter.

What remains is the research contribution, the base + complement split of milestone J2; an unencrypted disk, which will tell what BitLocker costs on a real model; Stone's biomimetic study, the weight genome; and the gate that will give this engine back to relais. The thesis exists in two complete versions, English and French; every figure is regenerated from the committed measurement data.
\end{otherlanguage}

\chapter*{Résumé}
\addcontentsline{toc}{chapter}{Résumé}
\markboth{Résumé}{Résumé}
\begin{otherlanguage}{french}
Le Tome~I de \emph{troncons} avait mesuré un portable ordinaire (16\,Go de RAM, une RTX\,4060 de 8\,Go, un disque NVMe chiffré), démontré le toit du problème et décidé, par le calcul, de construire d'abord une architecture exacte : servir par la carte graphique un modèle bien plus gros que toute la mémoire de la machine, ses experts compressés sans perte sur le disque et décodés à l'instant où un calcul en a besoin, dans une fenêtre de mémoire qui ne dépend pas de la taille du modèle. Ce second tome est celui du système : le programme, construit en Rust avec des noyaux CUDA par une équipe de cinq lots aux interfaces gelées, et son premier jalon, jugé sur gpt-oss-120b (63,4\,Go, quatre fois la RAM, huit fois la mémoire de la carte) contre douze critères déclarés avant la première ligne de code.

Le tome donne l'architecture et la structure du programme ; le format \texttt{.trc} au bit près, un tronçon par expert et des flux rangés par niveau de raffinement ; le chemin des octets, du disque à la mémoire de la carte, par la lecture sans tampon et la mémoire épinglée ; les noyaux, dont le décodeur rANS à table par voie, avec la preuve de leur exactitude et de leur déterminisme ; une référence en double précision et la théorie du juge, fondée sur la divergence de Kullback et Leibler ; le cache d'experts, l'ordonnanceur et le moteur, avec leurs lemmes.

Les résultats tiennent en quatre nombres, chacun recalculé par un second chemin depuis les passes commitées. \emph{Mémoire} : 384\,Mo de pointe pour le modèle de 63,4\,Go, la même à 3,7\,Mo près que pour un modèle cinq fois plus petit, quand llama.cpp en tient 10,9\,Go (11,4\,Go dans sa configuration la plus rapide) ; la mémoire ne suit plus la taille du modèle. \emph{Exactitude} : des logits identiques au bit entre le fichier codé et le fichier brut sur 5\,120 pas, d'une construction et d'un ordre de calcul à l'autre. \emph{Vitesse} : 1,965 jeton par seconde avec la garde d'intégrité, 2,22 fois la meilleure configuration de llama.cpp, et un préremplissage sept fois plus rapide ; deux rapports antérieurs, $\times 3{,}85$ puis $\times 2{,}79$, sont barrés parce que la base était mal réglée (la balayer l'a portée de 0,57 à 0,88 jeton/s). \emph{Justesse} : plus près des mathématiques exactes du modèle que llama.cpp (KL 0,0031 contre 0,0607 sur le 120b), et, avec un cache KV en f32, égal à elles au format près : tout l'écart venait de l'arrondi f16 du cache.

Un juge qui n'avait écrit ni la spécification ni le code a refait les chiffres : son décodeur, écrit depuis la seule spécification, redonne les tronçons au bit près ; sa propre passe avant en double précision confirme la référence à l'arrondi du format près. Il a trouvé deux défauts graves (une passe qui se fige sans message, un processus enfant orphelin qui bloque son pilote) et cinq moyens, dont une corruption d'octets bruts invisible jusque dans le texte. La ronde de corrections les a corrigés ou mesurés en une matinée, chacun avec un témoin qui tombe sur l'ancien binaire ; une garde de production rend désormais toute corruption d'un tronçon lu en erreur fatale nommée, pour 1,6\,\% de vitesse à disque égal. Dix critères tiennent sur douze ; les deux autres attendent une décision de l'arbitre.

Restent la contribution de recherche, la découpe base + complément du jalon J2 ; un disque non chiffré, qui dira ce que BitLocker coûte sur un vrai modèle ; l'étude biomimétique de Stone, le génome des poids ; et la porte qui rendra ce moteur à relais. Le mémoire existe en deux versions complètes, anglaise et française ; chaque figure se régénère depuis les données de mesure commitées.
\end{otherlanguage}
).

\chapter*{Abstract}
\addcontentsline{toc}{chapter}{Abstract}
\markboth{Abstract}{Abstract}
\begin{otherlanguage}{english}
Volume~I of \emph{troncons} measured an ordinary laptop (16\,GB of RAM, an 8\,GB RTX\,4060, an encrypted NVMe disk), proved the roof of the problem and decided, by calculation, to build an exact architecture first: serve through the graphics card a model much larger than all the memory of the machine, its experts compressed losslessly on disk and decoded at the moment a computation needs them, within a memory window that does not depend on the size of the model. This second volume is the volume of the system: the program, built in Rust with CUDA kernels by a team of five lots with frozen interfaces, and its first milestone, judged on gpt-oss-120b (63.4\,GB, four times the RAM, eight times the card's memory) against twelve criteria declared before the first line of code.

The volume gives the architecture and structure of the program; the \texttt{.trc} format to the bit, one chunk per expert and streams stored by refinement level; the byte path, from the disk to the card's memory, through unbuffered reads and pinned memory; the kernels, among them the rANS decoder with per-lane tables, with the proof of their exactness and determinism; a double-precision reference and the theory of the judge, founded on the Kullback--Leibler divergence; the expert cache, the scheduler and the engine, with their lemmas.

The results fit in four numbers, each recomputed by a second path from the committed passes. \emph{Memory}: a 384\,MB peak for the 63.4\,GB model, the same to within 3.7\,MB as for a model five times smaller, where llama.cpp holds 10.9\,GB (11.4\,GB in its fastest configuration); memory no longer follows the size of the model. \emph{Exactness}: logits identical to the bit between the coded and the raw file over 5\,120 steps, from one build and one order of computation to another. \emph{Speed}: 1.965 tokens per second with the integrity guard, 2.22 times the best configuration of llama.cpp, and a prefill seven times faster; two earlier ratios, $\times 3.85$ then $\times 2.79$, are struck out because the baseline was badly tuned (sweeping it raised it from 0.57 to 0.88 token/s). \emph{Correctness}: closer to the exact mathematics of the model than llama.cpp (KL 0.0031 against 0.0607 on the 120b), and, with an f32 KV cache, equal to them up to the format: the whole gap came from the f16 rounding of the cache.

A judge who wrote neither the specification nor the code redid the numbers: his decoder, written from the specification alone, reproduces the chunks to the bit; his own double-precision forward pass confirms the reference up to the rounding of the format. He found two serious defects (a pass that hangs without a message, an orphan child process that blocks its driver) and five medium ones, among them a corruption of raw bytes invisible down to the text. The round of corrections corrected or measured them in one morning, each with a control that fails on the old binary; a production guard now turns any corruption of a chunk read into a named fatal error, for 1.6\,\% of speed at equal disk. Ten criteria out of twelve hold; the other two await a decision of the arbiter.

What remains is the research contribution, the base + complement split of milestone J2; an unencrypted disk, which will tell what BitLocker costs on a real model; Stone's biomimetic study, the weight genome; and the gate that will give this engine back to relais. The thesis exists in two complete versions, English and French; every figure is regenerated from the committed measurement data.
\end{otherlanguage}

\chapter*{Résumé}
\addcontentsline{toc}{chapter}{Résumé}
\markboth{Résumé}{Résumé}
\begin{otherlanguage}{french}
Le Tome~I de \emph{troncons} avait mesuré un portable ordinaire (16\,Go de RAM, une RTX\,4060 de 8\,Go, un disque NVMe chiffré), démontré le toit du problème et décidé, par le calcul, de construire d'abord une architecture exacte : servir par la carte graphique un modèle bien plus gros que toute la mémoire de la machine, ses experts compressés sans perte sur le disque et décodés à l'instant où un calcul en a besoin, dans une fenêtre de mémoire qui ne dépend pas de la taille du modèle. Ce second tome est celui du système : le programme, construit en Rust avec des noyaux CUDA par une équipe de cinq lots aux interfaces gelées, et son premier jalon, jugé sur gpt-oss-120b (63,4\,Go, quatre fois la RAM, huit fois la mémoire de la carte) contre douze critères déclarés avant la première ligne de code.

Le tome donne l'architecture et la structure du programme ; le format \texttt{.trc} au bit près, un tronçon par expert et des flux rangés par niveau de raffinement ; le chemin des octets, du disque à la mémoire de la carte, par la lecture sans tampon et la mémoire épinglée ; les noyaux, dont le décodeur rANS à table par voie, avec la preuve de leur exactitude et de leur déterminisme ; une référence en double précision et la théorie du juge, fondée sur la divergence de Kullback et Leibler ; le cache d'experts, l'ordonnanceur et le moteur, avec leurs lemmes.

Les résultats tiennent en quatre nombres, chacun recalculé par un second chemin depuis les passes commitées. \emph{Mémoire} : 384\,Mo de pointe pour le modèle de 63,4\,Go, la même à 3,7\,Mo près que pour un modèle cinq fois plus petit, quand llama.cpp en tient 10,9\,Go (11,4\,Go dans sa configuration la plus rapide) ; la mémoire ne suit plus la taille du modèle. \emph{Exactitude} : des logits identiques au bit entre le fichier codé et le fichier brut sur 5\,120 pas, d'une construction et d'un ordre de calcul à l'autre. \emph{Vitesse} : 1,965 jeton par seconde avec la garde d'intégrité, 2,22 fois la meilleure configuration de llama.cpp, et un préremplissage sept fois plus rapide ; deux rapports antérieurs, $\times 3{,}85$ puis $\times 2{,}79$, sont barrés parce que la base était mal réglée (la balayer l'a portée de 0,57 à 0,88 jeton/s). \emph{Justesse} : plus près des mathématiques exactes du modèle que llama.cpp (KL 0,0031 contre 0,0607 sur le 120b), et, avec un cache KV en f32, égal à elles au format près : tout l'écart venait de l'arrondi f16 du cache.

Un juge qui n'avait écrit ni la spécification ni le code a refait les chiffres : son décodeur, écrit depuis la seule spécification, redonne les tronçons au bit près ; sa propre passe avant en double précision confirme la référence à l'arrondi du format près. Il a trouvé deux défauts graves (une passe qui se fige sans message, un processus enfant orphelin qui bloque son pilote) et cinq moyens, dont une corruption d'octets bruts invisible jusque dans le texte. La ronde de corrections les a corrigés ou mesurés en une matinée, chacun avec un témoin qui tombe sur l'ancien binaire ; une garde de production rend désormais toute corruption d'un tronçon lu en erreur fatale nommée, pour 1,6\,\% de vitesse à disque égal. Dix critères tiennent sur douze ; les deux autres attendent une décision de l'arbitre.

Restent la contribution de recherche, la découpe base + complément du jalon J2 ; un disque non chiffré, qui dira ce que BitLocker coûte sur un vrai modèle ; l'étude biomimétique de Stone, le génome des poids ; et la porte qui rendra ce moteur à relais. Le mémoire existe en deux versions complètes, anglaise et française ; chaque figure se régénère depuis les données de mesure commitées.
\end{otherlanguage}
).

\chapter*{Abstract}
\addcontentsline{toc}{chapter}{Abstract}
\markboth{Abstract}{Abstract}
\begin{otherlanguage}{english}
Volume~I of \emph{troncons} measured an ordinary laptop (16\,GB of RAM, an 8\,GB RTX\,4060, an encrypted NVMe disk), proved the roof of the problem and decided, by calculation, to build an exact architecture first: serve through the graphics card a model much larger than all the memory of the machine, its experts compressed losslessly on disk and decoded at the moment a computation needs them, within a memory window that does not depend on the size of the model. This second volume is the volume of the system: the program, built in Rust with CUDA kernels by a team of five lots with frozen interfaces, and its first milestone, judged on gpt-oss-120b (63.4\,GB, four times the RAM, eight times the card's memory) against twelve criteria declared before the first line of code.

The volume gives the architecture and structure of the program; the \texttt{.trc} format to the bit, one chunk per expert and streams stored by refinement level; the byte path, from the disk to the card's memory, through unbuffered reads and pinned memory; the kernels, among them the rANS decoder with per-lane tables, with the proof of their exactness and determinism; a double-precision reference and the theory of the judge, founded on the Kullback--Leibler divergence; the expert cache, the scheduler and the engine, with their lemmas.

The results fit in four numbers, each recomputed by a second path from the committed passes. \emph{Memory}: a 384\,MB peak for the 63.4\,GB model, the same to within 3.7\,MB as for a model five times smaller, where llama.cpp holds 10.9\,GB (11.4\,GB in its fastest configuration); memory no longer follows the size of the model. \emph{Exactness}: logits identical to the bit between the coded and the raw file over 5\,120 steps, from one build and one order of computation to another. \emph{Speed}: 1.965 tokens per second with the integrity guard, 2.22 times the best configuration of llama.cpp, and a prefill seven times faster; two earlier ratios, $\times 3.85$ then $\times 2.79$, are struck out because the baseline was badly tuned (sweeping it raised it from 0.57 to 0.88 token/s). \emph{Correctness}: closer to the exact mathematics of the model than llama.cpp (KL 0.0031 against 0.0607 on the 120b), and, with an f32 KV cache, equal to them up to the format: the whole gap came from the f16 rounding of the cache.

A judge who wrote neither the specification nor the code redid the numbers: his decoder, written from the specification alone, reproduces the chunks to the bit; his own double-precision forward pass confirms the reference up to the rounding of the format. He found two serious defects (a pass that hangs without a message, an orphan child process that blocks its driver) and five medium ones, among them a corruption of raw bytes invisible down to the text. The round of corrections corrected or measured them in one morning, each with a control that fails on the old binary; a production guard now turns any corruption of a chunk read into a named fatal error, for 1.6\,\% of speed at equal disk. Ten criteria out of twelve hold; the other two await a decision of the arbiter.

What remains is the research contribution, the base + complement split of milestone J2; an unencrypted disk, which will tell what BitLocker costs on a real model; Stone's biomimetic study, the weight genome; and the gate that will give this engine back to relais. The thesis exists in two complete versions, English and French; every figure is regenerated from the committed measurement data.
\end{otherlanguage}

\chapter*{Résumé}
\addcontentsline{toc}{chapter}{Résumé}
\markboth{Résumé}{Résumé}
\begin{otherlanguage}{french}
Le Tome~I de \emph{troncons} avait mesuré un portable ordinaire (16\,Go de RAM, une RTX\,4060 de 8\,Go, un disque NVMe chiffré), démontré le toit du problème et décidé, par le calcul, de construire d'abord une architecture exacte : servir par la carte graphique un modèle bien plus gros que toute la mémoire de la machine, ses experts compressés sans perte sur le disque et décodés à l'instant où un calcul en a besoin, dans une fenêtre de mémoire qui ne dépend pas de la taille du modèle. Ce second tome est celui du système : le programme, construit en Rust avec des noyaux CUDA par une équipe de cinq lots aux interfaces gelées, et son premier jalon, jugé sur gpt-oss-120b (63,4\,Go, quatre fois la RAM, huit fois la mémoire de la carte) contre douze critères déclarés avant la première ligne de code.

Le tome donne l'architecture et la structure du programme ; le format \texttt{.trc} au bit près, un tronçon par expert et des flux rangés par niveau de raffinement ; le chemin des octets, du disque à la mémoire de la carte, par la lecture sans tampon et la mémoire épinglée ; les noyaux, dont le décodeur rANS à table par voie, avec la preuve de leur exactitude et de leur déterminisme ; une référence en double précision et la théorie du juge, fondée sur la divergence de Kullback et Leibler ; le cache d'experts, l'ordonnanceur et le moteur, avec leurs lemmes.

Les résultats tiennent en quatre nombres, chacun recalculé par un second chemin depuis les passes commitées. \emph{Mémoire} : 384\,Mo de pointe pour le modèle de 63,4\,Go, la même à 3,7\,Mo près que pour un modèle cinq fois plus petit, quand llama.cpp en tient 10,9\,Go (11,4\,Go dans sa configuration la plus rapide) ; la mémoire ne suit plus la taille du modèle. \emph{Exactitude} : des logits identiques au bit entre le fichier codé et le fichier brut sur 5\,120 pas, d'une construction et d'un ordre de calcul à l'autre. \emph{Vitesse} : 1,965 jeton par seconde avec la garde d'intégrité, 2,22 fois la meilleure configuration de llama.cpp, et un préremplissage sept fois plus rapide ; deux rapports antérieurs, $\times 3{,}85$ puis $\times 2{,}79$, sont barrés parce que la base était mal réglée (la balayer l'a portée de 0,57 à 0,88 jeton/s). \emph{Justesse} : plus près des mathématiques exactes du modèle que llama.cpp (KL 0,0031 contre 0,0607 sur le 120b), et, avec un cache KV en f32, égal à elles au format près : tout l'écart venait de l'arrondi f16 du cache.

Un juge qui n'avait écrit ni la spécification ni le code a refait les chiffres : son décodeur, écrit depuis la seule spécification, redonne les tronçons au bit près ; sa propre passe avant en double précision confirme la référence à l'arrondi du format près. Il a trouvé deux défauts graves (une passe qui se fige sans message, un processus enfant orphelin qui bloque son pilote) et cinq moyens, dont une corruption d'octets bruts invisible jusque dans le texte. La ronde de corrections les a corrigés ou mesurés en une matinée, chacun avec un témoin qui tombe sur l'ancien binaire ; une garde de production rend désormais toute corruption d'un tronçon lu en erreur fatale nommée, pour 1,6\,\% de vitesse à disque égal. Dix critères tiennent sur douze ; les deux autres attendent une décision de l'arbitre.

Restent la contribution de recherche, la découpe base + complément du jalon J2 ; un disque non chiffré, qui dira ce que BitLocker coûte sur un vrai modèle ; l'étude biomimétique de Stone, le génome des poids ; et la porte qui rendra ce moteur à relais. Le mémoire existe en deux versions complètes, anglaise et française ; chaque figure se régénère depuis les données de mesure commitées.
\end{otherlanguage}
